% Part: propositional-logic
% Chapter: propositional-logic
% Section: preliminaries

\documentclass[../../../include/open-logic-section]{subfiles}

\begin{document}

\olfileid{pl}{syn}{pre}

\olsection{Prasyarat}

\begin{thm}[\emph{Asas induksi ing !!{formula}s}]
  \ollabel{thm:induction}  
  Yen sawijining sipat~$P$ lumaku kanggo kabeh !!{formula}s atomik lan
  nduweni kahanan kaya mangkene
  \begin{enumerate}
    \tagitem{prvNot}{sipat iku lumaku kanggo $\lnot !A$ saben sipat iku
      lumaku kanggo~$!A$;}{}
    \tagitem{prvAnd}{sipat iku lumaku kanggo $(!A \land !B)$
      saben sipat iku lumaku kanggo $!A$ lan~$!B$;}{}
    \tagitem{prvOr}{sipat iku lumaku kanggo $(!A \lor !B)$
      saben sipat iku lumaku kanggo $!A$ lan~$!B$;}{}
    \tagitem{prvIf}{sipat iku lumaku kanggo $(!A \lif !B)$
      saben sipat iku lumaku kanggo $!A$ lan~$!B$;}{}
    \tagitem{prvIff}{sipat iku lumaku kanggo $(!A \liff !B)$
      saben sipat iku lumaku kanggo $!A$ lan~$!B$;}{}
  \end{enumerate}
  mula $P$ lumaku kanggo kabeh !!{formula}s.
\end{thm}

\begin{proof}
  $S$ dadi koleksi kabeh !!{formula}s kang nduweni sipat~$P$. Cetha yen
  $S \subseteq \Frm[L_0]$. $S$~nyukupi kabeh syarat ing
  \olref[fml]{defn:formulas}: himpunan iki ngemot kabeh !!{formula}s atomik lan
  katutup marang !!{operator}s. $\Frm[L_0]$ iku kelas paling cilik kang
  kaya mangkono, mula $\Frm[L_0] \subseteq S$. Dadi $\Frm[L_0] = S$, lan
  saben formula nduweni sipat~$P$.
\end{proof}

\begin{prop}\ollabel{prop:balanced}
   Saben !!{formula} ing~$\Frm[L_0]$ iku \emph{imbang}, yaiku duwe
   tandha kurung kiwa lan tengen kang padha cacahe.
\end{prop}

\begin{prob}
Buktekna \olref[pl][syn][pre]{prop:balanced}
\end{prob}

\begin{prop} \ollabel{prop:noinit}
Ora ana perangan wiwitan sejati saka !!a{formula} kang uga dadi
!!a{formula}.
\end{prop}

\begin{prob}
  Buktekna \olref[pl][syn][pre]{prop:noinit}
\end{prob}

\begin{prop}[Bisa Kawaca kanthi Unik]
Saben !!{formula}~$!A$ ing $ \Frm[L_0]$ duwe persis siji uraian
minangka salah siji wujud ing ngisor iki
\begin{enumerate}
\tagitem{prvFalse}{$\lfalse$.}{}

\tagitem{prvTrue}{$\ltrue$.}{}

\item $\Obj p_n$ kanggo sawijining $\Obj p_n \in  \PVar$.
  
\tagitem{prvNot}{$\lnot !B$ kanggo sawijining !!{formula}~$!B$.}{}

\tagitem{prvAnd}{$(!B \land !C)$ kanggo sawijining !!{formula}s $!B$ lan~$!C$.}{}

\tagitem{prvOr}{$(!B \lor !C)$ kanggo sawijining !!{formula}s $!B$ lan~$!C$.}{}

\tagitem{prvIf}{$(!B \lif !C)$ kanggo sawijining !!{formula}s $!B$ lan~$!C$.}{}

\tagitem{prvIff}{$(!B \liff !C)$ kanggo sawijining !!{formula}s $!B$ lan~$!C$.}{}
\end{enumerate}
Kajaba iku, uraian iki \emph{unik}.
\end{prop}

\begin{proof}
Kanthi induksi ing $!A$. Tuladhane, upamakna $!A$ duwe rong wacan kang
beda minangka $(!B \lif !C)$ lan $(!B' \lif !C')$. Banjur $!B$ lan
$!B'$ kudu padha (yen ora, salah sijine bakal dadi perangan wiwitan
sejati saka liyane); mula yen rong wacan $!A$ iku beda, panyebabe
mesthi $!C$ lan $!C'$ dadi rong wacan kang beda saka rerangken simbol
kang padha. Iki mokal miturut hipotesis induktif.
\end{proof}


\begin{defn}[Substitusi Seragam]
Yen $!A$ lan $!B$ iku !!{formula}s, lan $\Obj p_i$ iku !!{propositional
variable}, mula $\Subst{!A}{!B}{\Obj p_i}$ nglambangake asil ngganti
saben kedadeyan $\Obj p_i$ nganggo sawijining kedadeyan $!B$ ing $!A$;
semono uga, substitusi bebarengan saka $\Obj p_1$, \dots,~$\Obj p_n$ dening
!!{formula}s $!B_1$, \dots,~$!B_n$ dilambangake kanthi
$\SSubst{!A}{\subst{!B_1}{\Obj p_1},\dots,\subst{!B_n}{\Obj p_n}}$.
\end{defn}

\begin{prob} Kanggo saben limang !!{formula}s ing ngisor iki, temtokna apa
 !!{formula} kasebut bisa ditulis minangka substitusi \( \Subst{!A}{!B}{\Obj p_i} \)
 nalika \( !A \) iku (i) \( \Obj p_0 \); (ii) \( ( \lnot \Obj p_0 \land \Obj
 p_1) \); lan (iii) \( ( ( \lnot \Obj p_0 \lif \Obj p_1 ) \land \Obj
 p_2 ) \). Ing saben kasus, temtokna substitusi kang cocog.
  \begin{enumerate}
    \item \( \Obj p_1 \) 
    \item \( ( \lnot \Obj p_0 \land \Obj p_0 ) \)
    \item \( ( ( \Obj p_0 \lor \Obj p_1 ) \land \Obj p_2 )  \)
    \item \( \lnot ( ( \Obj p_0 \lif \Obj p_1 ) \land \Obj p_2 ) \)
    \item \( (( \lnot ( \Obj p_0 \lif \Obj p_1 ) \lif ( \Obj p_0 \lor \Obj p_1 )) \land \lnot ( \Obj p_0 \land \Obj p_1 )) \)
  \end{enumerate}
\end{prob}

\begin{prob}
Wenehana definisi $\Subst{!A}{!B}{p}$ kang kaku sacara matematis
kanthi induksi.
\end{prob}

\end{document}
